\documentclass[11pt,a4paper]{article}
\usepackage{gtu-el-yazisi}
\begin{document}
\notbilgi{Alan Teorisi}{Gauss--Ostrogradski teoremi, Gauss yasası}{Ders8.pdf}

\sayfa{1}
\section{Gauss--Ostrogradski Teo:}

\elcizimi[0.45]{sekiller/Ders8-s1-1.png}[1]

\[
\int_{v}\mathrm{div}\,\vec A=\int_{S}\vec A\cdot ds
\]

\elcizimi[0.4]{sekiller/Ders8-s1-2.png}[1]

\[
\vec E=\frac{q}{4\pi\varepsilon_0}\,\frac{\vec e_r}{r^2}=\frac{q\cdot\vec r}{4\pi\varepsilon_0 r^3}
\]
\[
\vec D\triangleq\varepsilon_0\vec E=\frac{q}{4\pi}\,\frac{\vec r}{r^3}
\]
$\swarrow$ Deplasman \\ (yer değiştirme vektörü)

\[
\int_{S}\vec D\cdot ds=\int_{v}\mathrm{div}\,\vec D\cdot dv=\int_{v}\mathrm{div}\left(\frac{q}{4\pi}\,\frac{\vec r}{r^3}\right)dv
\]
\[
=\frac{q}{4\pi}\int_{v}{\color{kirmizikalem}-}\,\mathrm{div}\left({\color{kirmizikalem}-}\frac{\vec r}{r^3}\right)dv
\]
\[
=-\frac{q}{4\pi}\int_{v}\underbrace{\mathrm{div}\,\mathrm{grad}\left(\frac{1}{r}\right)}_{\Delta(1/r)}dv
\]

{\color{morkalem}
\[
\frac{\partial}{\partial x}\vec e_x+\frac{\partial}{\partial y}\vec e_y+\frac{\partial}{\partial z}\vec e_z
\]
\[
f(x,y,z)\ \Rightarrow\ \mathrm{grad}\,f=\vec\nabla\cdot f=\frac{\partial f}{\partial x}\vec e_x+\frac{\partial f}{\partial y}\vec e_y+\frac{\partial f}{\partial z}\vec e_z
\]}

\sayfa{2}
{\color{morkalem}
\[
\mathrm{div}\,\mathrm{grad}\,f=\frac{\partial^2 f}{\partial x^2}+\frac{\partial^2 f}{\partial y^2}+\frac{\partial^2 f}{\partial z^2}=\Delta f
\qquad\underline{\text{laplacian}}
\]}

\[
r=\sqrt{x^2+y^2+z^2}\ \Rightarrow\ \frac{1}{r}=(x^2+y^2+z^2)^{-1/2}
\]
\[
\Delta\Bigl((x^2+y^2+z^2)^{-1/2}\Bigr)
\]

veya küresel koordinatlarda çözersiniz

\begin{align*}
f(r,\theta,\phi)\footnotemark&=\frac{1}{r^2}\,\frac{\partial}{\partial r}\Bigl(r^2\,\frac{\partial f}{\partial r}\Bigr)+\underbrace{\frac{1}{r^2\sin\theta}\,\frac{\partial}{\partial\theta}\Bigl(\sin\theta\,\frac{\partial f}{\partial\theta}\Bigr)}_{\textcolor{kirmizikalem}{0}}\\
&\quad+\underbrace{\frac{1}{r^2\sin^2\theta}\,\frac{\partial^2 f}{\partial\phi^2}}_{\textcolor{kirmizikalem}{0}}
\end{align*}
\footnotetext{Claude notu: muhtemelen $\Delta f(r,\theta,\phi)$ olmalı.}

\[
f=\frac{1}{r}=f(r)
\]
\[
\frac{\partial(1/r)}{\partial\theta}=\frac{\partial(1/r)}{\partial\phi}=0
\]

\[
\Delta(1/r)=\frac{1}{r^2}\,\frac{\partial}{\partial r}\Bigl(r^2\,\underbrace{\frac{\partial}{\partial r}(1/r)}_{-1/r^2}\Bigr)
=\frac{1}{r^2}\,\frac{\partial}{\partial r}\Bigl[-r^2\,\frac{1}{r^2}\Bigr]
\]
\[
=-1\,\frac{\partial}{\partial r}(1)=0\ \Rightarrow\ (r\neq 0)
\]

$(r\neq 0)$ ise $q\notin v$ ise $\displaystyle\int_{S}\vec D\cdot ds=0$

\sayfa{3}
\elcizimi[0.4]{sekiller/Ders8-s3-1.png}[3]

{\color{kirmizikalem}$S'$: $r'$ yarıçaplı bir küre yüzeyi}

$v''$ bölgesi $S\cup S'$ yüzeylerini kuşatır

\[
\vec r\in v''\ \Rightarrow\ \Delta(1/r)=0
\]
\[
\Rightarrow\ \int_{v''}\mathrm{div}\,\vec D\cdot dv''=\int_{S+S'}\vec D\,(d\vec s+d\vec s')
\]
\[
\Rightarrow\ -\frac{q}{4\pi}\int_{v''}\overbrace{\Delta(1/r)}^{0}dv''
\]
\[
=\int_{S+S'}\vec D\,(d\vec s+d\vec s')=0\;=\int_{S}\vec D\,d\vec s+\int_{S'}\vec D\cdot d\vec s'=0
\]

\[
\int_{S}\vec D\cdot d\vec s=\,??\qquad(r=0\in v)\qquad q\in v
\]

\[
\int_{S}\vec D\cdot d\vec s=-\int_{S'}\vec D\cdot d\vec s'
\]
$q$'nun $S'$ de yarattığı $\vec D$: \quad $\vec D=\dfrac{q}{4\pi}\,\dfrac{\vec r'}{(r')^3}$

\elcizimi[0.2]{sekiller/Ders8-s3-2.png}[3]
\[
\begin{aligned}
ds&=a^2\sin\theta\,d\theta\,d\phi\\
dv&=r^2\sin\theta\,d\theta\,d\phi\footnotemark
\end{aligned}
\]
\footnotetext{Claude notu: muhtemelen $dv=r^2\sin\theta\,dr\,d\theta\,d\phi$ olmalı.}

\[
d\vec s'=-\vec e_{r'}\cdot\underbrace{r'^2\sin\theta\,d\theta\,d\phi}_{ds'}
\]
\[
\vec D\cdot d\vec s'=\frac{q}{4\pi}\,\frac{r'\cdot}{r'^3}\Bigl(\vec e_{r'}\cdot r'^2\sin\theta'\,d\theta\,d\phi\Bigr)
\]

\sayfa{4}
\[
\int_{S'}\vec D\cdot ds'=-\int_{\theta'=0}^{\pi}\int_{\phi'=0}^{2\pi}-\frac{q}{4\pi}\,\sin\theta'\,d\theta'\,d\phi'
\]
\[
=\frac{q}{4\pi}\,\underbrace{\int_{\theta'=0}^{\pi}\sin\theta'\,d\theta'}_{\textcolor{kirmizikalem}{2}}\ \underbrace{\int_{\phi'=0}^{2\pi}d\phi'}_{\textcolor{kirmizikalem}{2\pi}}
\]
\[
=\frac{q}{4\pi}\cdot 4\pi=q
\]

\[
\int_{S}\vec D\cdot d\vec s=Q_{\text{iç}}+\frac{Q_{\text{yüzey}}}{2}
\]

\elcizimi[0.5]{sekiller/Ders8-s4-1.png}[4]

$1/2$ yarıçaplı küre yüzeyi\\
tekillik için koyulur yüzey üzerinde yük olduğu için

\[
\int_{S}\vec D\cdot d\vec s=-2C+\left(\frac{-4C}{2}\right)
\]
\[
=-2C-2C=-4C
\]

\[
[\vec D]=\mathrm{C/m^2}
\]
\[
[\vec E]=\mathrm{V/m}\qquad\qquad [V]=\mathrm{V_o}
\]
\[
\int_{S}\vec D\cdot d\vec s=-2C-4C=-6C
\]

\[
\int_{S'}\vec D\cdot d\vec s'=-2C-4C+1C=-5C
\]

\sayfa{5}
{\color{kirmizikalem}Sınav 2 ya da 3 soru}

\[
\int_{S}\vec D\cdot d\vec s=\int_{v}\mathrm{div}\,\vec D\cdot dv=\underbrace{\int_{v}\rho\,dv}_{v\text{'deki }Q\text{ yükü}}
\]
\[
=Q_{\text{iç}}
\]
\[
\rho=\frac{dQ}{dv}\qquad(v\text{ içindeki h.y.y.}\quad[\rho]=\mathrm{C/m^3})
\]

\formulkutu[Maxwell]{\mathrm{div}\,\vec D=\rho}

Boşlukta
\[
\vec D=\varepsilon_0\vec E\ \Rightarrow\ \mathrm{div}\,\vec E=\frac{\rho}{\varepsilon_0}
\]
\[
\vec E=-\mathrm{grad}\,V
\]
\[
\Rightarrow\ -\underbrace{\mathrm{div}\,\mathrm{grad}\,V}_{\Delta V}=\frac{\rho}{\varepsilon_0}\ \Rightarrow\ \Delta V=-\frac{\rho}{\varepsilon_0}
\]

$v$'de $\rho\neq 0\ \Rightarrow\ v$'deki $V$ pot fonksiyonun
\[
\kutula{\Delta V=-\rho/\varepsilon_0}\ \Rightarrow\ V=\frac{1}{4\pi\varepsilon_0}\int_{v}\frac{\rho}{r'}\,dv
\]
$\downarrow$ Poisson denklemi

$v$'de $\rho=0$ (yani yük yok) $\Rightarrow$
\formulkutu[Laplace]{\Delta V=0}

\elcizimi[0.75]{sekiller/Ders8-s5-1.png}[5]

\sayfa{6}
\noindent S: 32\\ pr 5

\elcizimi[0.3]{sekiller/Ders8-s6-1.png}[6]

\[
v=\frac{4\pi a^3}{3}
\]
homojen yayılmış $\swarrow$
\[
\kutula{\rho\triangleq\frac{dQ}{dv}}\ \Rightarrow\ \frac{Q}{v}=\rho=\frac{3Q}{4\pi a^3}\qquad(\mathrm{C/m^2})\footnotemark
\]
\footnotetext{Claude notu: hacimsel yük yoğunluğu için birim C/m$^3$ olmalı.}

\[
Q=\int_{v}\rho\,dv=\int_{r=0}^{a}\int_{\theta=0}^{\pi}\int_{\phi=0}^{2\pi}\frac{3Q}{4\pi a^3}\,r^2\sin\theta\,d\theta\,d\phi\,dr
\]
\[
=\frac{3Q}{4\pi a^3}\int_{r=0}^{a}r^2\,dr\int_{\theta=0}^{\pi}\sin\theta\,d\theta\int_{\phi=0}^{2\pi}d\phi
\]
\[
=\frac{3Q}{4\pi a^3}\,\frac{a^3}{3}\cdot 2\cdot 2\pi=Q
\]

\[
V=V(r)\qquad,\ r>a\qquad\vec E=-\mathrm{grad}\,V=-\frac{\partial V}{\partial r}\,\vec e_r
\]
\[
\vec D=\varepsilon_0\vec E=-\underbrace{\varepsilon_0\,\frac{\partial V}{\partial r}}_{D(r)}\,\vec e_r
\]
\[
d\vec s=r^2\sin\theta\,d\theta\,d\phi\,\vec e_r
\]
\[
\int_{S_r}\vec D\cdot d\vec s=Q\ \Rightarrow\ \int_{S_r}\underbrace{D(r)\,\vec e_r}_{\vec D}\ \overset{1}{\cdot}\ \underbrace{\vec e_r\cdot r^2\sin\theta\,d\theta\,d\phi}_{d\vec s_r}=Q
\]
\[
\int_{\theta=0}^{\pi}\int_{\phi=0}^{2\pi}D(r)\,r^2\sin\theta\,d\theta\,d\phi=Q
\]

\sayfa{7}
\[
\Rightarrow\ D(r)\,r^2\int_{\theta=0}^{\pi}\sin\theta\,d\theta\int_{\phi=0}^{2\pi}d\phi=Q
\]
\[
D(r)\,r^2\,4\pi=Q
\]
\[
D(r)=\frac{Q}{4\pi r^2}\ \Rightarrow\ \vec D=\frac{Q}{4\pi r^2}\,\vec e_r
\]

o zaman merkezde tek bir yük gibi küre davranır $r>a$ için

\[
\vec D(r)=\frac{Q}{4\pi r^2}\,\vec e_r=-\varepsilon_0\,\frac{\partial V}{\partial r}\,\vec e_r
\]
\[
\textcolor{turuncukalem}{V=\int\frac{\partial V}{\partial r}\,dr=\int\frac{-Q}{4\pi\varepsilon_0 r^2}\,dr}\ \Rightarrow\ V(r)=\frac{Q}{4\pi\varepsilon_0}\,\frac{1}{r}+V_0
\]
\[
r\to\infty\ \Rightarrow\ V=0\ \Rightarrow\ V_0=0
\]

\[
V(r)=\frac{Q}{4\pi\varepsilon_0}\,\frac{1}{r}+V_0
\]
\[
\lim_{r\to\infty}V(r)=V_0=0
\]
\[
r>a:\qquad V=\frac{Q}{4\pi\varepsilon_0}\cdot\frac{1}{r}
\]
\[
\Delta V=\frac{Q}{4\pi\varepsilon_0}\,\underbrace{\Delta(1/r)}_{0}=0
\]

\textcolor{kirmizikalem}{A}\quad\underline{$0<r<a$}

\elcizimi[0.5]{sekiller/Ders8-s7-1.png}[7]

\sayfa{8}
\[
\vec D=D(r)\,\vec e_r
\]
\[
\int_{S_r}\vec D\cdot d\vec s_r=\int_{\theta=0}^{\pi}\int_{\phi=0}^{2\pi}D(r)\,\underbrace{\vec e_r\cdot\vec e_r}_{1}\,r^2\sin\theta\,d\theta\,d\phi=\int_{v_r}\rho\,dv_r
\]
\[
D(r)\,r^2\int_{\theta=0}^{\pi}\sin\theta\,d\theta\int_{0}^{2\pi}d\phi=\int\frac{3Q}{4\pi a^3}\,dv_r
\]
\[
\Rightarrow\ D(r)\,r^2\,4\pi=\frac{3Q}{4\pi a^3}\,v_r=\frac{3Q}{4\pi a^3}\,\frac{4\pi r^3}{3}
\]

\[
\int_{v_r}dv_r=v_r=\frac{4\pi r^3}{3}
\]
\[
\int_{v_r}dv_r=\int_{0}^{r}\int_{\theta=0}^{\pi}\int_{\phi=0}^{2\pi}\underbrace{r^2\sin\theta\,d\theta\,d\phi\,dr}_{dv_r=dr\,ds_r}
\]
\[
=\int_{r=0}^{r}r^2\,dr\underbrace{\int_{\theta=0}^{\pi}\sin\theta\,d\theta}_{2}\underbrace{\int_{\phi=0}^{2\pi}d\phi}_{2\pi}
\]
\[
\frac{r^3}{3}\Big|_0^r=\frac{r^3}{3}
\]
\[
\frac{4\pi r^3}{3}
\]

\[
\Rightarrow\ D(r)\,r^2\,4\pi=\frac{Q\cdot r^3}{a^3}
\]
\[
\Rightarrow\ D(r)=\frac{Q\,r}{4\pi a^3}
\]
\[
\vec D=\kutula{\frac{Q}{4\pi a^3}\,r}\,\vec e_r=\kutula{-\varepsilon_0\,\frac{\partial V}{\partial r}}\,\vec e_r
\]

\sayfa{9}
\[
\frac{\partial V}{\partial r}=-\frac{Q}{4\pi\varepsilon_0 a^3}\,r
\]
\[
V=-\frac{Q}{8\pi\varepsilon_0 a^3}\,r^2+\widetilde{V_0}
\]
\[
\underbrace{V\Big|_{r=a^+}}_{\downarrow}=\underbrace{V\Big|_{r=a^-}}_{\downarrow}
\]
\[
\frac{Q}{4\pi\varepsilon_0}\,\frac{1}{r}\Big|_{r=a}=-\frac{Q}{8\pi\varepsilon_0 a^3}\,r^2\Big|_{r=a}+\widetilde{V_0}
\]
\[
\frac{Q}{4\pi\varepsilon_0}\,\frac{1}{a}=-\frac{Q}{8\pi\varepsilon_0 a}+\widetilde{V_0}
\]
\[
\widetilde{V_0}=\frac{3Q}{8\pi\varepsilon_0 a}
\]

$0<r<a$

\formulkutu{V(r)=-\frac{Q}{8\pi\varepsilon_0 a}\left[\frac{r^2}{a^2}-3\right]}

$r>a$:
\[
V(r)=\frac{Q}{4\pi\varepsilon_0}\,\frac{1}{r}
\]

\elcizimi[0.4]{sekiller/Ders8-s9-1.png}[9]

\[
V(r)=-\frac{Q\,r^2}{8\pi\varepsilon_0 a^3}+\frac{3Q}{8\pi\varepsilon_0 a}
\]
\[
\Delta V=\frac{1}{r^2}\frac{\partial}{\partial r}\left(r^2\frac{\partial V}{\partial r}\right)=-\frac{Q}{8\pi\varepsilon_0 a^3}\,\Delta(r^2)
\]

\sayfa{10}
\[
=-\frac{Q}{8\pi\varepsilon_0 a^3}\cdot\frac{1}{r^2}\cdot\frac{\partial}{\partial r}\Big[\overbrace{r^2\,\frac{\partial(r^2)}{\partial r}}^{2r^3}\Big]
\]
\[
=-\frac{Q}{4\pi\varepsilon_0 a^3}\cdot\frac{3\,r^2}{r^2}\footnotemark=-\kutula{\frac{3Q}{4\pi a^3}}\;\frac{\rho}{\varepsilon_0}
\]
\footnotetext{El yazısında pay ve paydadaki $r^2$ terimleri çizilerek sadeleştirilmiş.}
\[
=-\frac{\rho}{\varepsilon_0}
\]

\end{document}
